% TOP_PROBLEM: {"id": 7500111, "title": "Strength of the Carlson–Simpson Lemma", "category_id": 18, "category": {"id": 18, "name": "logic", "display_name": "Logic", "description": "Problems in mathematical logic, model theory, proof theory, and finite model theory.", "slug": "logic", "order_index": 18, "created_at": "2026-07-31T15:26:25.671Z"}, "status": "partially_solved"}
% FIRST_POSTED: null
% ENTRY_KIND: statement_only
% Imported from: batch20.tex; UnsolvedMath ID 7500111
\documentclass[11pt]{article}
\usepackage[a4paper,margin=25mm]{geometry}
\usepackage{amsmath,amssymb,mathtools}
\usepackage{fontspec,unicode-math}
\setmainfont{Latin Modern Roman}
\setsansfont{Latin Modern Sans}
\setmathfont{Latin Modern Math}
\usepackage{xurl,hyperref}
\hypersetup{colorlinks=true,linkcolor=blue,urlcolor=blue,pdftitle={UnsolvedMath Problem 7500111}}
\newcommand{\sref}[2]{\hyperlink{source:#1:#2}{[S#2]}}
\newcommand{\cataloguescope}[1]{\paragraph{Mathematical question.}#1}
\begin{document}
% UnsolvedMath ID 7500111; source catalogue code AMR-074-0111
\setcounter{section}{7500110}
\section{Strength of the Carlson–Simpson Lemma}\label{7500111}
{\small\sffamily UnsolvedMath ID 7500111 · Logic}
\subsection{Definitions and mathematical statement}

\paragraph{Shared notation.}$\mathbb N=\{1,2,\ldots\}$, and $\mathbb Z,\mathbb Q,\mathbb R,\mathbb C$ denote the integers, rationals, reals and complex numbers. Any other conventions are specified locally.


Work in second-order arithmetic over $\mathsf{RCA}_0$, the base theory with $\Delta^0_1$ comprehension and $\Sigma^0_1$ induction. The strength of a principle means which subsystems prove it and which subsystems it implies over that base; equivalence requires both implications. Let $A$ be a finite nonempty alphabet and $\mathbb N_0=\{0,1,2,\ldots\}$. An infinite variable word $W$ is a sequence in $A\cup\{x_i:i\in\mathbb N_0\}$ in which every variable occurs and all occurrences of $x_i$ precede every occurrence of $x_{i+1}$. For $\bar a=a_0\cdots a_{k-1}\in A^{<\omega}$, $W(\bar a)$ is formed by replacing $x_i$ by $a_i$ for $i<k$ and stopping just before the first $x_k$. The Carlson--Simpson principle $\mathsf{CS}$ says that every finite colouring of $A^{<\omega}$ admits such a $W$ with $\{W(\bar a):\bar a\in A^{<\omega}\}$ monochromatic.
\paragraph{Statement recorded in the supplied source.}
Determine the reverse-mathematical strength of $\mathsf{CS}$. \sref{7500111}{2}
\subsection{Short English statement}
Which set-existence axioms are needed to find an infinite variable word whose finite substitutions all have the same colour?
\subsection{Sources}
\begin{description}
\item[\hypertarget{source:7500111:1}{S1}] ULAM.AI and the UnsolvedMath Contributors, UnsolvedMath: A Curated Collection of Open Mathematics Problems, record 7500111 (AMR-074-0111). CC BY 4.0. \url{https://huggingface.co/datasets/ulamai/UnsolvedMath}.
\item[\hypertarget{source:7500111:2}{S2}] Antonio Montalbán, Open Questions in Reverse Mathematics, author preprint dated 17 August 2010, Question 11, printed p. 5. \url{https://math.berkeley.edu/~antonio/papers/questionsRM.pdf}.
\end{description}
\label{end:7500111}
\end{document}
